% !TEX root = ../../main.tex
% ============================================================================
% Standard Landscape; Platonic Introduction.
% ============================================================================

\providecommand{\cA}{\mathcal{A}}
\providecommand{\cH}{\mathcal{H}}
\providecommand{\cM}{\mathcal{M}}
\providecommand{\cW}{\mathcal{W}}
\providecommand{\Walg}{\mathcal{W}}
\providecommand{\fg}{\mathfrak{g}}
\providecommand{\fsl}{\mathfrak{sl}}
\providecommand{\Vir}{\mathrm{Vir}}
\providecommand{\GRT}{\mathrm{GRT}_1}
\providecommand{\Q}{\mathbb{Q}}
\providecommand{\Z}{\mathbb{Z}}
\providecommand{\ChirAlg}{\mathrm{ChirAlg}}
\providecommand{\Einf}{E_\infty}
\providecommand{\Eone}{E_1}
\providecommand{\Etwo}{E_2}
\providecommand{\MC}{\mathrm{MC}}

\chapter{The Five-Archetype Landscape as a moduli atlas}
\label{ch:part-iii-platonic-introduction}

Part~\ref{part:bar-cobar-engine} produced the bar coalgebra $B(A_b)$
for any augmented chart algebra. Part~\ref{part:bulk} produced its
derived chiral centre $Z^{\mathrm{der}}_{\mathrm{ch}}(A_b)$. The
present part asks: across the standard landscape of vertex algebras,
how does the chart $A_b$ stratify, and what coordinates classify the
strata? The answer is the $5{\times}5$ $\kappa$-stratification matrix
across the five archetypes
$\mathsf G/\mathsf L/\mathsf C/\mathsf M/\mathsf B$.

The classical reading of the Kac--Wakimoto table prints one scalar per
vertex algebra. Heisenberg gets $c=1$; free fermion gets $c=1/2$;
affine~$V_k(\fsl_2)$ gets $c = 3k/(k+2)$; Virasoro gets~$c$;
$\cW_N$ principal gets a function of $c$ and~$N$; the $\beta\gamma$
system gets $c = 2(6\lambda^2 - 6\lambda + 1)$. A reader confronted
with this list sees a parameter taking values in~$\mathbb{C}$, and is
invited to imagine that the landscape is indexed by that parameter.

The reading is at once true and incomplete. Heisenberg and free
fermion both have $c=1$ at $\lambda = 0$ for the fermion; they are
not the same algebra. Virasoro at $c=1/2$ and the Ising model
minimal model both have central charge~$1/2$; the first is a universal
algebra, the second a simple quotient, and their shadow invariants
distinguish them. Affine $V_k(\fsl_2)$ at $k = -1$ has $c = -3$; so
does a rescaled Heisenberg doublet; the bar coalgebras differ in
every degree beyond the second. The central charge is a coarse
projection.

The classification is canonical only after the scalar is replaced by
a vector of typed projections. The five-slot fingerprint
\begin{equation}
 \varphi(\cA) \;=\; \bigl(p_{\max},\; r_{\max},\;
 \chi_{\mathrm{VOA}},\; n_{\mathrm{strong}},\; \mathrm{coset}\bigr)
 \label{eq:part-iii-fingerprint}
\end{equation}
records the finite row data from which the central charge, the
classical collision residue, the scalar conductor, and the shadow
class are read on the standard rows. In the language of
Theorem~\ref{thm:fingerprint-completeness}
\textup{(}Chapter~\ref{ch:three-invariants}\textup{)}, $\varphi$ is
a fingerprint projection of the modular-Koszul package. The full
chiral bar coalgebra requires the OPE, ordered coalgebra, and
Hochschild/centre data in addition to~$\varphi$. The landscape is the
fingerprint atlas of~$\cM_{\ChirAlg}^{\mathrm{std}}$: it records
finite invariants and coarse shadow strata while the full OPE, the bar
coalgebra, and the derived centre remain separate data.

This introduction to Part~\ref{part:standard-landscape}
names that atlas, catalogues the charts, and records the theorems that
make the atlas rigorous. The eight sections of
Part~\ref{part:standard-landscape} correspond to strata. The
associator action enters only on ordered-chain presentations equipped
with a $\GRT(\Q)$-torsor, and the historical faces $F_1$--$F_7$ are
comparison charts only after their carriers and transition maps have
been fixed. Every family enters through its typed bar, centre, scalar,
and coset data.

\section{Deficiency of the central-charge classification}
\label{sec:part-iii-deficiency}

Four symptoms show that $c$ alone cannot sort the standard landscape.

\begin{enumerate}[leftmargin=*]
 \item \textbf{Degeneracy in~$c$.} Distinct chiral algebras share a
 central charge. The Heisenberg algebra~$\cH_1$ at level~$1$ has
 $c = 1$; the rank-$1$ lattice vertex operator algebra
 $V_\Z$ has~$c = 1$; the simple $(2,3)$ Virasoro minimal model has
 $c = -22/5$, equal to neither. The fingerprint coordinates
 $(p_{\max}, r_{\max}, n_{\mathrm{strong}}, \mathrm{coset})$ separate
 these three.

 \item \textbf{No distinction of class.} The central charge does not
 read the shadow class. Virasoro at generic~$c$ is class~$M$ with
 $S_4 = 10/[c(5c+22)] \ne 0$ (Zamolodchikov 1985); the critical-level affine
 $V_{-h^\vee}(\fg)$ has $\kappa_{\mathrm{ch}} = 0$ and lives in a
 companion class~$\mathrm{FF}$
 \textup{(}Theorem~\ref{thm:fifth-class-FF}\textup{)} where the
 shadow metric degenerates. The bar cohomology of the first grows
 polynomially; the bar cohomology of the second is infinite-dimensional
 by the Feigin--Frenkel center. The scalar~$c$ detects neither
 dichotomy.

 \item \textbf{No Drinfeld-Sokolov transport.} The reduction
 $V_k(\fg) \mapsto H^\bullet_{\mathrm{DS}}(V_k(\fg),\, f_{\mathrm{prin}})
 = \cW_{\varkappa}(\fg, f_{\mathrm{prin}})$
 takes a class~$L$ algebra to a class~$M$ algebra, preserves the
 Koszul conductor (Theorem~\ref{thm:DS-fingerprint-transport}), and
 changes $c$ by a non-trivial function. The fingerprint transports
 slot-by-slot; the central charge loses the slot structure.

 \item \textbf{No duality.} Bar-cobar inversion
 $\Omega\bar B^{\mathrm{ch}}(\cA)\to\cA$ and the Verdier/Koszul
 partner~$\cA^!$ are different constructions. On the named Verdier
 lanes the scalar sums are family-dependent: the Virasoro
 matter-ghost central-charge sum is~$26$, the operadic scalar
 complementarity is recorded by $K^\kappa$, and the
 Bershadsky--Polyakov row has its own value. A single scalar cannot
 carry these typed complementarity laws; the fingerprint's coset slot
 records which Verdier lane is being used.
\end{enumerate}

Each symptom forces a slot of~$\varphi$. The OPE pole
order~$p_{\max}$ separates Heisenberg ($p_{\max}=2$) from the free
fermion ($p_{\max}=1$) and from Virasoro ($p_{\max}=4$). The shadow
depth~$r_{\max}$ cuts the non-degenerate locus into four archetype
classes~$G,L,C,M$. The Hodge-filtered
supertrace~$\chi_{\mathrm{VOA}}$ reads parity: the
symplectic boson (class~$G$) and the symplectic fermion (class~$C$)
share OPE pole order but differ in parity. The strong-generator
count~$n_{\mathrm{strong}}$ detects the number of
BRST-independent fields. The coset slot records the pair
$(\mathrm{Com}(\cH_{\kappa}, \cA), \cA/\mathrm{Com})$ at
non-critical level and its degenerate
$\mathrm{coset}_{\mathrm{FF}}$ at critical level.

The classification is then the one the chiral bar demands. Part~\ref{part:standard-landscape}
is the catalogue of charts on the moduli stack
$\cM_{\ChirAlg}^{\mathrm{std}}$: the objects are standard
chiral-algebra packages, the fingerprint quotient records their
finite row data, and the central-charge reading is the first
coordinate projection of the atlas, truncated at the trace.

\section{Eight sections of the standard landscape}
\label{sec:part-iii-eight-sections}

Part~\ref{part:standard-landscape} carries eight charts on
$\cM_{\ChirAlg}^{\mathrm{std}}$.

\begin{description}
 \item[\S III.0 Platonic introduction.]
 \label{sec:part-iii-0}
 Central-charge deficiency. The eight-section atlas and the five
 Platonic theorems replace scalar classification by the full
 fingerprint.

 \item[\S III.1 Free-field atoms.]
 \label{sec:part-iii-1}
 Chapters~\ref{chap:free-fields} and~\ref{chap:beta-gamma}: Heisenberg
 $\cH_\kappa$, free fermion~$\psi$, $bc$ ghosts at weight~$\lambda$,
 $\beta\gamma$ system, lattice VOAs $V_L$. Free-field OPE; all exactly
 Koszul. The atoms at~$p_{\max} \in \{1,2\}$ and
 $r_{\max} \in \{2,3,4\}$, with the symplectic boson-vs-fermion
 dichotomy at~$\chi_{\mathrm{VOA}} = \pm 1$ made explicit.

 \item[\S III.2 Affine Kac--Moody.]
 \label{sec:part-iii-2}
 Chapter~\ref{chap:kac-moody}: $V_k(\fg)$ for every simple
 Lie algebra~$\fg$, with the
 level-parameter stratification $k \in \mathbb{C}$ decomposed into
 admissible, generic, and critical strata. At generic~$k$ the algebra
 is class~$L$; at critical $k = -h^\vee$ the shadow metric
 degenerates and the algebra enters class~$\mathrm{FF}$; at
 admissible levels the Feigin--Frenkel involution
 $k \leftrightarrow -k - 2h^\vee$ fixes the Koszul dual.

 \item[\S III.3 Virasoro and $\cW$-algebras.]
 \label{sec:part-iii-3}
 Chapter~\ref{chap:w-algebras}: the stress-tensor hierarchy
 $\mathrm{Vir}_c \hookrightarrow \cW_N \hookrightarrow \cW_\infty[\mu]$
 via Drinfeld--Sokolov reduction. All class~$M$ at generic
 level, with the shadow coefficient
 $S_4(\cW_N)$ controlled by the universal formula
 $\kappa(\cW_N) = c\,(H_N - 1)$ where
 $H_N = \sum_{j=1}^{N} 1/j$. The quartic
 Virasoro pole forces the infinite shadow tower; the
 $\cW_N$ tower refines it by~$N-1$ additional stress tensors.

 \item[\S III.4 Super lineage.]
 \label{sec:part-iii-4}
 Chapter~\ref{chap:n2-sca} and the super-Virasoro entries
 of the census: $N=1, N=2, N=4$ superconformal algebras, symplectic
 fermion, super-Yangian. The fingerprint acquires a
 $\Z/2$-grading on the coset slot; the super-variant
 $\Gsuper = \GRT(\Q) \rtimes (\Z/2)^{|\text{odd generators}|}$ acts.
 Symplectic boson and symplectic fermion share~$p_{\max}$ and
 $r_{\max}$ but differ in~$\chi_{\mathrm{VOA}}$.

 \item[\S III.5 Exceptional and sporadic.]
 \label{sec:part-iii-5}
 Chapter~\ref{ch:bershadsky-polyakov}, Chapter~\ref{ch:y-algebras},
 the Monster and Leech lattice moonshine chapters. Bershadsky--Polyakov
 as subregular DS reduction (hook-type, self-dual with Koszul
 conductor $K = 196$); Gaiotto--Rapcak $Y_{N_1, N_2, N_3}[\psi]$
 with $S_3$-triality permuting the coset slot; Monster $V^\natural$
 as $\Z/2$-orbifold of the Leech lattice VOA
 $V_{\Lambda_{24}}^+$. These families sit at distinguished
 $\mathbb{C}$-points
 where two or more stratum walls meet.

 \item[\S III.6 Coset constructions.]
 \label{sec:part-iii-6}
 Chapter~\ref{ch:logarithmic-w-algebras} and the symmetric orbifold
 chapter: GKO coset,
 Kac--Roan--Wakimoto hook-type reductions, the logarithmic triplet
 $\cW(p)$. The coset slot of~$\varphi$ becomes the governing datum;
 the classification of cosets is an orbit under the relative Weyl
 group of the coset pair.

 \item[\S III.7 Landscape as moduli stack.]
 \label{sec:part-iii-7}
 Chapter~\ref{ch:landscape-census}: the census, reread as the atlas
 of charts on $\cM_{\ChirAlg}^{\mathrm{std}}$. The five
 strata~$G, L, C, M, \mathrm{FF}$ each correspond to a locally closed
 substack; the Koszul involution
 $\cA \mapsto \cA^!$ is an involution of the atlas exchanging
 pairs of charts; each family in Table~\ref{tab:master-invariants}
 is a $\mathbb{C}$-point with an explicit fingerprint.

 \item[\S III.8 Exhaustiveness theorem.]
 \label{sec:part-iii-8}
 The exhaustiveness clause: the finite standard realizability locus
 in $\mathrm{Fing}$ is exactly the image of the standard-landscape
 charts under~$\varphi$. The proof is the fibre product of the five
 theorem carriers below.
\end{description}

The standard form carries eight charts on a single moduli stack,
with the Koszul involution, the Drinfeld--Sokolov transport, and the
associator comparison action visible as structure on the atlas.

\section{Five Platonic theorems}
\label{sec:part-iii-five-theorems}

The eight-section atlas is held together by five theorem surfaces.
Each has a stated comparison carrier; each is a
statement about~$\cM_{\ChirAlg}^{\mathrm{std}}$ whose proof uses the
chapters in Parts~I--II and the fingerprint machinery of
Chapter~\ref{ch:three-invariants}.

\subsection{Theorem A: landscape as fingerprint atlas}
\label{sec:part-iii-theorem-A}

\begin{theorem}[Standard landscape as a fingerprint atlas;
 \ClaimStatusConditional]
\label{thm:part-iii-landscape-as-moduli-stack}
Let $\cM_{\ChirAlg}^{\mathrm{std}}$ denote the moduli stack of
standard-landscape chiral algebras on the fixed Koszul/PBW finite
windows used in Part~\ref{part:standard-landscape}. The canonical
fingerprint map
\[
 \varphi:\; \cM_{\ChirAlg}^{\mathrm{std}} \;\longrightarrow\;
 \mathrm{Fing} \;=\;
 \mathbb{Z}_{\ge 0} \times (\{2,3,4,\infty\} \cup \{\mathrm{FF}\})
 \times \Z \times \mathbb{Z}_{\ge 0} \times \mathrm{Coset},
\]
stratifies $\cM_{\ChirAlg}^{\mathrm{std}}$ into the five finite-row
classes $G,L,C,M,\mathrm{FF}$. Every chiral algebra in the standard
landscape gives a point of this stack with a specified fingerprint,
and the fingerprint quotient records precisely the finite row data
listed in the census. It does not identify two chiral algebras with
the same fingerprint unless their full OPE and bar-coalgebra data
are separately identified.
\end{theorem}
\begin{proof}
The stack keeps the full chiral algebra as the object and attaches
the five-slot fingerprint as a projection. The five archetype classes
partition the image of~$\varphi$ by
Theorem~\ref{thm:five-class-stratum}. Every family in
Table~\ref{tab:master-invariants} has an explicit fingerprint in the
census, so each gives a point of
$\cM_{\ChirAlg}^{\mathrm{std}}$. The non-injectivity statement is
Theorem~\ref{thm:fingerprint-completeness}: the fingerprint is read
from the modular-Koszul package but does not reconstruct the full
package.
\end{proof}

\subsection{Theorem B: central-charge functional}
\label{sec:part-iii-theorem-B}

\begin{theorem}[The central charge is a functional of the fingerprint;
 \ClaimStatusConditional]
\label{thm:part-iii-central-charge-is-emergent}
Let $\cA$ be a chiral algebra in the standard landscape with
fingerprint $\varphi(\cA) = (p_{\max}, r_{\max}, \chi_{\mathrm{VOA}},
n_{\mathrm{strong}}, \mathrm{coset})$ at
$\kappa_{\mathrm{ch}}(\cA) \ne 0$. The central charge~$c(\cA)$ is
recovered from~$\varphi$ by a family-dependent formula whose
parameters are read off from $\varphi$:
\begin{equation}
 c(\cA) \;=\; \Psi_{\mathrm{coset}}\bigl(
 p_{\max}, r_{\max}, \chi_{\mathrm{VOA}},
 n_{\mathrm{strong}}\bigr)
 \label{eq:part-iii-central-charge-functional}
\end{equation}
where $\Psi_{\mathrm{coset}}$ is the central-charge functional
\textup{(}Definition~\ref{def:central-charge-functional},
Remark~\ref{rem:central-charge-functional-examples}\textup{)} attached
to the coset slot. In particular, the central charge is an
$\MC$-coinvariant of the universal Maurer--Cartan element
$\Theta_\cA$ and coincides with twice the Koszul invariant
$\kappa_{\mathrm{ch}}(\cA)$ for the Virasoro stratum.
\end{theorem}
\begin{proof}
The central charge is read off from the degree-$2$ truncation of the
shadow tower: $c(\cA) = 2\kappa_{\mathrm{ch}}(\cA)$ in the Virasoro
stratum, and in general by the family-dependent Sugawara shift of
Theorem~\ref{thm:genus-universality}. The fingerprint slots determine
the shift: $p_{\max}$ selects the pole-order hypothesis of the
Sugawara construction; the coset slot encodes the central extension
class needed to convert $\kappa_{\mathrm{ch}}$ to~$c$; the remaining
slots fix the numerical coefficients. Explicit formulae:
\begin{itemize}[leftmargin=*]
 \item Virasoro stratum: $c(\mathrm{Vir}_c) = 2\kappa_{\mathrm{ch}}$.
 \item Affine $V_k(\fg)$ (class~$L$): $c = k \dim\fg / (k + h^\vee)$,
 equivalently $c = 2 \dim\fg\cdot(k + h^\vee)/ (2h^\vee) - \dim\fg
 = 2\kappa_{\mathrm{ch}} - \dim\fg$ absorbing the dimension shift
 via~$n_{\mathrm{strong}}$.
 \item $\cW_N$ principal: $c = 2\kappa_{\mathrm{ch}} / (H_N - 1)$
 \textup{(}inverting
 $\kappa_{\mathrm{ch}}(\cW_N) = c\,(H_N - 1)$\textup{)}.
 \item Heisenberg~$\cH_\kappa$: $c = 1$ is read from the strong
 generator count~$n_{\mathrm{strong}} = 1$ directly.
 \item $\beta\gamma$ / $bc$ system at weight~$\lambda$: the coset slot
 pins $\lambda$, and the functional
 $\Psi_{\mathrm{coset}}$ reads $c = 2(6\lambda^2 - 6\lambda + 1)$ or
 $c = 1 - 3(2\lambda - 1)^2$ according to
 parity~$\chi_{\mathrm{VOA}}$.
\end{itemize}
In every case the formula is single-valued on the stratum and
determined by~$\varphi(\cA)$. Thus $c$ is a coset-dependent
functional of the five fingerprint slots on the stated rows.
\end{proof}

\subsection{Theorem C: r-matrix from PBW and OPE}
\label{sec:part-iii-theorem-C}

\begin{theorem}[The classical r-matrix is determined by PBW and
 OPE data; \ClaimStatusConditional]
\label{thm:part-iii-rmatrix-emergent-from-pbw}
Let $\cA$ be a chiral algebra in the standard landscape with a chosen
set of strong generators indexed by~$n_{\mathrm{strong}}$ and a PBW
basis on the augmentation ideal. The classical r-matrix
\begin{equation}
 r(z) \;=\; \Res^{\mathrm{coll}}_{0,2}(\Theta_\cA)
 \;\in\; \cA^! \otimes \cA^! [\![z^{-1}]\!]
\end{equation}
is determined by the pair (PBW basis, OPE). Explicitly, for each
stratum the formula is:
\begin{enumerate}[label=\textup{(\roman*)}, leftmargin=*]
 \item Heisenberg (class~$G$): $r(z) = \kappa / z$, simple pole,
 level prefix~$\kappa$ mandatory (the trace-form convention in
 which $r(z)|_{\kappa=0}=0$ recovers the abelian limit).
 \item Affine $V_k(\fg)$ (class~$L$, trace-form convention):
 $r(z) = k\,\Omega_{\mathrm{tr}} / z$, where~$\Omega_{\mathrm{tr}}
 \in \fg \otimes \fg$ is the trace-form Casimir; $k = 0$ gives
 $r = 0$ consistent with the abelian limit.
 \item Virasoro (class~$M$): $r(z) = (c/2)/z^3 + 2T/z$, cubic and
 simple pole (the descendant coefficient $T(w)/(z-w)^2$ is absorbed
 by the $d\log$ kernel into the simple-pole term).
 \item $\cW_N$ (class~$M$): $r(z) = \sum_{n=2}^{N} r^{(n)}(z)$
 with $r^{(n)}(z) = \kappa^{(n)}/z^{2n-1} + \text{(descendant
 corrections)}$, one component per stress tensor.
 \item $\beta\gamma, bc$ (class~$C$): $r(z) = 0$ at the level of
 the collision residue; the non-trivial content enters at the
 quartic shadow~$S_4$.
\end{enumerate}
In every case the formula follows from the augmentation-ideal PBW
decomposition, the OPE pole structure, and the collision residue of
the universal MC element~$\Theta_\cA$. The r-matrix is therefore a
functional of~$(p_{\max}, n_{\mathrm{strong}}, \mathrm{coset})$; it
is not an independent datum.
\end{theorem}
\begin{proof}
The collision residue~$r(z) = \Res^{\mathrm{coll}}_{0,2}(\Theta_\cA)$
is computed directly from the degree-$2$ truncation of the shadow
tower \textup{(}Theorem~\ref{thm:mc2-bar-intrinsic}\textup{)}. The
PBW basis fixes the strong generators; the OPE pole structure
determines the coefficients of the collision residue by reading off
pole orders and Laurent coefficients. For each class:
\begin{itemize}[leftmargin=*]
 \item Class~$G$: OPE has only a double pole, collision residue is
 a simple pole in~$z$, coefficient is the level.
 \item Class~$L$: OPE has a double pole carrying the Casimir, plus a
 simple pole carrying the Lie bracket; d-log absorption yields a
 single-pole r-matrix with Casimir weight. Sugawara shift relates
 the classical r-matrix to~$\kappa_{\mathrm{ch}}$ via
 $\mathrm{av}(r) + \dim\fg/2 = \kappa_{\mathrm{ch}}$ for non-abelian
 affine.
 \item Class~$M$: quartic OPE pole forces a cubic plus simple pole
 collision residue; the cubic coefficient is~$c/2$, the simple
 coefficient is the stress tensor~$T$.
 \item Class~$C$: no positive binary collision pole remains after
 the $d\log$ absorption of the simple generator pole. The non-trivial
 shadow content enters through the quartic contact class, for
 $\beta\gamma$ with $S_4=-5/12$ and $r_{\max}=4$.
\end{itemize}
The pole-by-pole readings determine the collision residue. $\quad\square$
\end{proof}

\subsection{Theorem D: associator action on the fingerprint quotient}
\label{sec:part-iii-theorem-D}

\begin{theorem}[Associator action on the fingerprint quotient;
 \ClaimStatusConditional]
\label{thm:part-iii-grt-orbit-structure}
Let $\cM_{\ChirAlg}^{\mathrm{fp}}$ be the fingerprint quotient of the
completed standard-landscape packages for which the ordered
Fulton--MacPherson chain model is equipped with a Drinfeld associator
torsor and the face comparison maps
\[
\rho_i:T_i\longrightarrow \cF_{\mathrm{fp}}(\cA)
\]
have been fixed. Then $\GRT(\Q)$ acts on the associator choices,
preserves the fingerprint projection, and acts on the fibre-product
image of the comparison maps. Under these hypotheses:
\begin{enumerate}[label=\textup{(\roman*)}, leftmargin=*]
 \item each fingerprint stratum
 $\cM^{(\varphi_0),\mathrm{fp}}_{\ChirAlg}$ is stable under
 associator transport;
 \item the historical seven faces
 $F_1, F_2, \ldots, F_7$ of the classical r-matrix
 \textup{(}Vol~I bar hub, Dimofte--Niu--Py R-twist, Kniznik-Zamolodchikov
 PVA r-matrix, Gaiotto--Zeng commuting differentials, Drinfeld Yangian,
 Feigin--Frenkel--Reshetikhin Gaudin, class-$M$ top $A_\infty$
 $m_3$\textup{)} are $\Q$-rational comparison charts in the
 $\GRT(\Q)$-torsor structure on $\Face(\cA)$
 \textup{(}Vol~II Chapter~\ref{V2-ch:grt-seven-faces}\textup{)};
 \item two additional Q-rational faces exist:
 $F_8$ \textup{(}Brown motivic, MZV-coupled\textup{)} and
 $F_9$ \textup{(}Willwacher operadic, via
 $H^0(\mathsf{GC}_2) = \mathfrak{grt}_1$\textup{)};
 \item the super-variant
 $\Gsuper = \GRT(\Q) \rtimes (\Z/2)^{|\text{odd}|}$
 acts on the super-lineage chart $\S\,\ref{sec:part-iii-4}$,
 absorbing the symplectic boson-vs-fermion
 asymmetry into its~$\Z/2$-factor.
\end{enumerate}
\end{theorem}
\begin{proof}
The ordered FM-chain model is functorial in the associator. Changing
the associator by an element of $\GRT(\Q)$ transports the operations
and hence transports every carrier $T_i$ for which the comparison map
$\rho_i$ is defined. The fingerprint slots are read from pole order,
shadow depth, Hodge supertrace, strong-generator count, and coset
data; these are invariant under a change of associator presentation.
Thus the $\GRT(\Q)$-action preserves the fingerprint quotient and the
fibre-product image of the maps~$\rho_i$.

The seven historical faces are charts on one comparison object only
after the carriers and maps have been supplied. Vol~II constructs the
face torsor and the two additional motivic/operadic charts $F_8,F_9$;
this theorem records the induced action on the
Part~\ref{part:standard-landscape} fingerprint quotient. The
super-lineage chart carries a $\Z/2$-grading on the coset slot, which
gives the displayed semidirect extension.
\end{proof}

\subsection{Theorem E: atlas completeness}
\label{sec:part-iii-theorem-E}

\begin{theorem}[Standard realizability locus of the fingerprint;
 \ClaimStatusConditional]
\label{thm:part-iii-atlas-completeness}
Let $\mathrm{Fing}^{\mathrm{std}}\subset\mathrm{Fing}$ be the
smallest subset containing the fingerprints of the free-field,
affine, $\beta\gamma$/ghost, Virasoro/$\cW$, super, exceptional,
sporadic, coset, and critical-level companion charts, and closed under
the tensor, Drinfeld--Sokolov, coset, and Verdier operations whose
PBW/Koszul hypotheses are verified in the cited chapters. Then:
\begin{enumerate}[label=\textup{(\roman*)}, leftmargin=*]
 \item every $\varphi_0\in\mathrm{Fing}^{\mathrm{std}}$ is realized
 by at least one point
 $\cA\in\cM_{\ChirAlg}^{\mathrm{std}}(\mathbb{C})$;
 \item over such a point, the groupoid of ordered Fulton--MacPherson
 chain presentations equipped with a Drinfeld associator carries the
 $\GRT(\Q)$ torsor action of
 Theorem~\ref{thm:part-iii-grt-orbit-structure};
 \item on the perfect Verdier lanes, the Koszul partner
 $\cA \mapsto \cA^!$ sends $\varphi_0$ to the dual fingerprint
 vector~$\varphi_0^!$ under the complementarity involution
 \textup{(}Theorem~\ref{thm:central-charge-complementarity}
 for the central charge; the general fingerprint involution
 follows from the Koszul-equivariance of $\varphi$,
 Remark~\ref{rem:fingerprint-koszul-symmetry}\textup{)}.
\end{enumerate}
\end{theorem}
\begin{proof}
The definition of $\mathrm{Fing}^{\mathrm{std}}$ is constructive.
The free-field, affine, $\beta\gamma$/ghost, Virasoro/$\cW$,
super, exceptional, sporadic, and coset chapters provide the initial
witnesses. Tensor product, Drinfeld--Sokolov reduction, and coset
formation preserve the stated finite PBW/Koszul windows by
Theorem~\ref{thm:pole-depth-independence},
Theorem~\ref{thm:d-alg-r-max-bijection}, and
Theorem~\ref{thm:DS-fingerprint-transport}; the critical-level
degeneracies are placed in the companion class by
Theorem~\ref{thm:fifth-class-FF}. This proves realization for the
standard locus.

The associator statement is the torsor action on ordered
Fulton--MacPherson chain presentations of a fixed standard package,
transported through the comparison maps~$\rho_i$ of
Theorem~\ref{thm:part-iii-grt-orbit-structure}. On the perfect
Verdier lanes, Koszul-equivariance of~$\varphi$
\textup{(}Remark~\ref{rem:fingerprint-koszul-symmetry}\textup{)}
identifies the dual fingerprint. $\quad\square$
\end{proof}

\section{The landscape as moduli atlas}
\label{sec:part-iii-moduli-atlas}

Theorems A through E organise Part~\ref{part:standard-landscape} into an atlas on
$\cM_{\ChirAlg}^{\mathrm{std}}$. The atlas has the following shape.

\begin{itemize}[leftmargin=*]
 \item \textbf{Five strata.} $\cM_{\ChirAlg}^{\mathrm{std}}$ is stratified into
 five locally closed substacks:
 $G$ (free-field, $r_{\max} = 2$),
 $L$ (affine non-critical, $r_{\max} = 3$),
 $C$ ($\beta\gamma$/ghost, $r_{\max} = 4$),
 $M$ (Virasoro/$\cW_N$, $r_{\max} = \infty$),
 $\mathrm{FF}$ (critical-level companion,
 $\kappa_{\mathrm{ch}} = 0$).
 Each stratum is an open substack of the closure of a unique
 irreducible component.

 \item \textbf{Fourteen affine charts.} On the Koszulness moduli
 scheme $\cM_{\mathrm{Kosz}}(\cA)$
 \textup{(}Theorem~\ref{thm:koszulness-moduli-kp}\textup{)}, each
 family has a distinguished Drinfeld associator
 under which its Koszul-bar-cobar inversion is
 unconditional. The fourteen-chart atlas
 \textup{(}Table~\ref{tab:master-invariants-chart}\textup{)} records
 the map
 $\text{family} \mapsto \text{associator chart}$:
 class~$G$ sits at $\Phi_{\mathrm{KZ}}$;
 class~$L$ at $\Phi_{\mathrm{ell}}$ (Enriquez);
 class~$C$ at $\Phi_{\mathrm{AT}}$ (Alekseev--Torossian);
 class~$M$ at $\Phi_{\mathrm{Kon}}$ (Kontsevich integral);
 class~$\mathrm{FF}$ at $\Phi_{\mathrm{dR,B}}$ (Hodge--Betti).

 \item \textbf{Koszul involution.} $\cA \mapsto \cA^!$ acts on the
 atlas, pairing charts:
 $G \leftrightarrow G$ (self-paired for Heisenberg),
 $L \leftrightarrow L$ (via $k \leftrightarrow -k - 2h^\vee$),
 $C \leftrightarrow C$ ($\beta\gamma \leftrightarrow bc$),
 $M \leftrightarrow M$ (via $c \leftrightarrow 13 - c$ in the
 operadic convention, $c + c' = 26$ in the matter-ghost convention).

 \item \textbf{Drinfeld--Sokolov transport.} The DS reduction
 $V_k(\fg) \mapsto \cW_\varkappa(\fg, f)$ is a functorial transport
 $L \to M$ on the atlas, sending the level stratification of~$L$
 into the stress-tensor-count stratification of~$M$ by the
 shift~$\varkappa = -(k + h^\vee)^{-1}$
 \textup{(}Kac--Wakimoto\textup{)}. The fingerprint transports
 slot-by-slot under
 Theorem~\ref{thm:DS-fingerprint-transport}.

 \item \textbf{GRT action.} $\GRT(\Q)$ acts on each stratum
 preserving the fingerprint. The orbit
 structure exhibits seven (respectively nine, with $F_8$ and $F_9$)
 $\Q$-rational representatives per orbit, each corresponding to a
 historical face of the classical r-matrix.
\end{itemize}

\section{Paramodular K3 chart over $\overline{\mathcal A_2}$}
\label{sec:part-iii-paramodular-chart}

The Vol III comparison atlas attached to
$\cM_{\ChirAlg}^{\mathrm{std}}$ carries a distinguished chart
indexed by the Baily--Borel--Freitag stratification of
$\overline{\mathcal A_2}$. Write
$\mathfrak g_{\Delta_5} = \mathrm{Borch}(F_3,\phi^{K3}_{0,1})$ for
the K3-BKM Lie superalgebra from the hyperbolic lattice
$\Lambda_{\mathrm{Muk},II}^{2,1}$; its chiral bialgebra realisation
$\mathbf H_{\Delta_5}$ has Weyl denominator
$\Delta_5 \in S_5(K(1))$ (Gritsenko paramodular weight-$5$ cusp form).
Four sibling functors index the chart:
\begin{description}
 \item[$\Psi$] reflective interior, producing the generic
 $\mathbf H_{\Delta_5}$-points;
 \item[$\Psi^{\deg}$] Klingen cusp, producing the super-affine
 degeneration;
 \item[$\Psi^{\mathrm{tor}}$] Humbert divisor, producing the
 quantum-toroidal specialisation;
 \item[$\Psi^{\mathrm{metap}}$] metaplectic branch, producing
 Conway $V^{s\natural}$.
\end{description}
Completeness plus absence of a fifth stratum is inscribed in
Chapter~\ref{V3-ch:k3-chiral-bialgebra-platonic}: the Humbert
divisor coincides with the Argyres--Douglas locus of the
class-$\mathcal S$ SCFT $\mathcal T[A_1,\Sigma_{0,24}]$, where at
$E_{\tau_1}\times E_{\tau_2}$ the Seiberg--Witten curve degenerates
to $(A_1,A_{2N-1})$ (Gaiotto--Moore--Neitzke $2009$, Example~$8.3$).
The four paramodular siblings average to the four shadow archetypes
$G/L/C/M$ of Theorem~\ref{thm:part-iii-landscape-as-moduli-stack};
the $FF$ companion is the critical-level degeneration of the
reflective interior.

\begin{theorem}[K3 master $L$-value identity on the atlas;
 \ClaimStatusProvedElsewhere]
\label{thm:part-iii-k3-master-L-value}
The genus-$1$ free energy of $\mathbf H_{\Delta_5}$ on the
paramodular K3 chart is
\[
\log Z^{(1)}_{\mathbf H_{\Delta_5}}
\;=\;-\log\bigl\|\Delta_5\bigr\|^2_{\mathrm{Quillen}}
\;=\;-\log\Delta_5 \;-\; 24\,L'(0,\Delta_5,\mathrm{std}),
\qquad
24 = \chi_{\mathrm{top}}(K3),
\]
by Bruinier--K\"uhn $2003$ composed with Bismut--Gillet--Soul\'e $1988$.
The $L$-function is standard, not adjoint-spinor: Pitale--Saha--Schmidt
$2014$ CAP factorisation combined with the Waldspurger squaring rule
rule out $L'(0,\Delta_{10},\mathrm{ad}^0)$. The source theorem is
Chapter~\ref{V3-ch:k3-chiral-bialgebra-platonic}.
\end{theorem}

\begin{remark}[Bloch--Kato Selmer tangent to the paramodular chart]
\label{rem:part-iii-selmer-tangent}
The Bloch--Kato group
$H^1_f(\mathbb Q,\mathrm{std}\,\rho_{\Delta_5})$ is an arithmetic
comparison target for the paramodular chart at
$\mathbf H_{\Delta_5}$. Interpreting it as an actual tangent line to
the chiral-bialgebra chart requires a chiral trace map, compatibility
with the completed Hall source, and comparison with the Heegner
divisor class. Without these maps it is a source-recognition target,
not a constructed deformation of the fingerprint coset slot.
\end{remark}

\begin{remark}[Humbert--Heegner admissibility on the atlas]
\label{rem:part-iii-humbert-heegner-atlas}
The pentagon tower $\{\phi^{(n)}\}_{n\ge 3}$ on the K3 comparison
chart satisfies
$\phi^{(n)} = 0$ unless $n\equiv 3,5\pmod 8$, equivalently
$D_n = (n-3)/2\in\{0,1\}\pmod 4$. The polar-support cutoff for the
index-$1$ K3 Jacobi form concentrates the tower at $n=3$ and $n=5$:
the $H_3$ wall carries $\phi^{(3)}|_{H_3}$, and the one-term
$n=5$ coherence is
$\phi^{(5)}=2\cdot[\mathrm{gen}]^{\otimes 5}$ in the normalization of
Theorem~\ref{thm:bc-unconditional-depth-reduction}. For admissible
$n\ge11$ the polar coefficient vanishes.
\end{remark}

\section{Two filtrations on the standard landscape}
\label{sec:part-iii-reading-order}

Part~\ref{part:standard-landscape} carries two filtrations.
The atom filtration starts with the free-field chart
(\S~\ref{sec:part-iii-1}) and adds non-abelian OPE, stress tensors,
parity, exceptional dualities, cosets, the moduli chart, and
exhaustiveness. At each step one slot of~$\varphi$ becomes active.

The moduli filtration starts at the atlas
(\S~\ref{sec:part-iii-7}) and descends through
$G,L,C,M,\mathrm{FF}$. On each stratum the functional
$\Psi_{\mathrm{coset}}$ specializes, the r-matrix is read from the
PBW/OPE carrier, and the associator comparison action relates the
$\Q$-rational charts. The two filtrations meet where the
stratum-by-stratum content becomes the atlas of
$\cM_{\ChirAlg}^{\mathrm{std}}$.

\section{Comparison carriers}
\label{sec:part-iii-dependency-map}

The five Platonic theorems
(\ref{thm:part-iii-landscape-as-moduli-stack},
\ref{thm:part-iii-central-charge-is-emergent},
\ref{thm:part-iii-rmatrix-emergent-from-pbw},
\ref{thm:part-iii-grt-orbit-structure},
\ref{thm:part-iii-atlas-completeness}) have the following carriers in
Parts~\ref{part:bar-complex}--\ref{part:characteristic-datum} and
Chapter~\ref{ch:three-invariants}.

\begin{center}
\begin{tabular}{p{3.8cm}|p{6.2cm}|p{4.5cm}}
\hline
\textbf{Theorem} & \textbf{Standard-landscape carrier} &
\textbf{Proof input} \\
\hline
\ref{thm:part-iii-landscape-as-moduli-stack}
 (landscape as fingerprint atlas)
 & \S III.7 (\texttt{landscape\_census}),
 \S III.8 (completeness)
 & \ref{thm:fingerprint-completeness},
 \ref{thm:five-class-stratum},
 \ref{thm:fifth-class-FF} \\
\hline
\ref{thm:part-iii-central-charge-is-emergent}
 (central-charge functional)
 & \S III.1 (\texttt{free\_fields}),
 \S III.2 (\texttt{kac\_moody}),
 \S III.3 (\texttt{w\_algebras})
 & \ref{thm:genus-universality},
 \ref{thm:central-charge-complementarity},
 \ref{thm:mc2-bar-intrinsic} \\
\hline
\ref{thm:part-iii-rmatrix-emergent-from-pbw}
 (r-matrix from PBW and OPE)
 & \S III.1, \S III.2, \S III.3
 & \ref{thm:mc2-bar-intrinsic},
 \ref{thm:pbw-koszulness-criterion} \\
\hline
\ref{thm:part-iii-grt-orbit-structure}
 (associator action)
 & \S III.4 (\texttt{n2\_superconformal}),
 \S III.5 (\texttt{bershadsky\_polyakov}, \texttt{y\_algebras}),
 \S III.7
 & Vol II \ref{thm:grt-seven-faces-torsor},
 \ref{thm:koszulness-moduli-kp},
 \ref{thm:DS-fingerprint-transport} \\
\hline
\ref{thm:part-iii-atlas-completeness}
 (standard realizability)
 & \S III.8 (exhaustiveness)
 & \ref{thm:pole-depth-independence},
 \ref{thm:d-alg-r-max-bijection},
 \ref{thm:DS-fingerprint-transport},
 \ref{thm:fifth-class-FF} \\
\hline
\end{tabular}
\end{center}

\section{The landscape, restated}
\label{sec:part-iii-landscape-restated}

The standard landscape is the fingerprint quotient of the finite-row
chiral packages carried by
$\cM_{\ChirAlg}^{\mathrm{std}}$. Every family is a
$\mathbb{C}$-point of this moduli stack; its central charge,
classical r-matrix, and Koszul conductor are explicit functionals of
a five-slot fingerprint; its place in the Koszul involution is visible
on the perfect Verdier lanes; its comparison with the other families
is a transport law on the atlas.

\begin{proposition}[Standard landscape restated as moduli atlas;
 \ClaimStatusConditional]
\label{prop:part-iii-landscape-restated}
For the fingerprint quotient fixed above, the content of Part~\ref{part:standard-landscape} is the
triple
\begin{equation}
 \bigl(\cM_{\ChirAlg}^{\mathrm{std}},\; \varphi,\; \{\rho_i:T_i\to\cF_{\mathrm{fp}}\}\bigr),
 \label{eq:part-iii-landscape-triple}
\end{equation}
where $\cM_{\ChirAlg}^{\mathrm{std}}$ is the standard-landscape
moduli stack on the stated Koszul/PBW finite windows, $\varphi$ is
the canonical fingerprint projection, and the maps~$\rho_i$ are the
typed face comparisons. The classical Kac--Wakimoto census is the
list of standard points with their finite row data; the central
charge, r-matrix, and scalar conductor are read by the typed
functionals above.
\end{proposition}
\begin{proof}
Composition of
Theorem~\ref{thm:part-iii-landscape-as-moduli-stack} (stack
existence), Theorem~\ref{thm:part-iii-central-charge-is-emergent}
(central-charge functional),
Theorem~\ref{thm:part-iii-rmatrix-emergent-from-pbw}
(r-matrix functional),
Theorem~\ref{thm:part-iii-grt-orbit-structure}
($\GRT$ associator action), and
Theorem~\ref{thm:part-iii-atlas-completeness}
(exhaustiveness). $\quad\square$
\end{proof}

\medskip

The landscape is the fingerprint atlas of standard chiral-algebra
families. Its strata are controlled by shadow depth and coset data;
its Koszul involution is a Verdier-lane operation on typed packages;
its Drinfeld--Sokolov transport is a functorial map between charts.
The Kac--Wakimoto list supplies standard points. Central charge is a
coset-dependent functional of the fingerprint, and the classical
r-matrix is read from the PBW and OPE data.
